\chapter{综合实验与总结}
\label{chap:integration}

\section{实验目的}

本章是整本实验指导书的收束章节，目标是帮助读者把前面各章拆开的模块重新串成一条完整的 DBMS 执行链路。完成本章后，读者应能够从一组 SQL 出发，说明 DBMS\_C 如何完成建表、插入、查询、修改、删除、提交和空结果查询，并能把这些操作对应到解析、计划、记录、事务、缓冲区、日志和文件页式存储等模块。

本章不新增自动测试，也不扩展系统功能，而是用综合实验、图表和模块边界总结，把前面已经验证过的 CTest 与人工 Trace 观察组织成最终复盘材料。

\section{综合实验主线}

综合实验建议使用一组短小但覆盖面较完整的 SQL。它既包含更新类命令，也包含 SELECT 查询和最终提交，适合作为人工运行、Trace 观察和源码调用链复盘的统一样例。

\begin{figure}[htbp]
  \centering
  \resizebox{0.80\textwidth}{!}{%
  \begin{tikzpicture}[
    node distance=8mm and 12mm,
    every node/.style={align=center},
    core/.style={dblevel, text width=32mm, minimum height=9mm},
    mod/.style={dbnode, text width=34mm, minimum height=10mm},
    side/.style={dbsubnode, text width=30mm, minimum height=10mm}
  ]
    \node[core] (sql) {SQL 输入};
    \node[mod, below=of sql] (parse) {第8章 Parse\\QueryData / CommandData};

    \node[mod, below left=10mm and 16mm of parse] (query) {第9章 Plan / Scan\\SELECT 查询};
    \node[mod, below right=10mm and 16mm of parse] (update) {第10章 Update\\更新命令分发};
    \node[side, right=16mm of update] (trace) {第11章 Trace\\CLI 过程观察};

    \node[mod, below=20mm of parse] (meta) {第7章 Metadata\\系统目录};
    \node[mod, below=of meta] (record) {第6章 Record\\记录布局与表扫描};

    \node[mod, below left=10mm and 10mm of record] (tx) {第5章 Transaction\\提交、回滚、锁};
    \node[mod, below right=10mm and 10mm of record] (buffer) {第4章 Buffer\\pin / dirty / flush};

    \node[mod, below=19mm of record] (file) {第2章 File / Page\\文件块与页};
    \node[side, below left=6mm and 8mm of tx] (log) {第3章 Log\\日志追加与刷盘};
    \node[side, below right=6mm and 8mm of buffer] (lib) {第1章 Lib\\基础工具};

    \draw[dbarrow] (sql) -- (parse);
    \draw[dbarrow] (parse) -- (query);
    \draw[dbarrow] (parse) -- (update);
    \draw[dbarrow] (query) -- (meta);
    \draw[dbarrow] (update) -- (meta);
    \draw[dbarrow] (meta) -- (record);
    \draw[dbarrow] (record) -- (tx);
    \draw[dbarrow] (record) -- (buffer);
    \draw[dbarrow] (tx) -- (log);
    \draw[dbarrow] (tx) -- (buffer);
    \draw[dbarrow] (buffer) -- (file);
    \draw[dbarrow] (lib) -- (buffer);

    \draw[dbdashline] (trace) -- (parse);
    \draw[dbdashline] (trace) -- (update);
    \draw[dbdashline] (trace) -- (tx);

  \end{tikzpicture}}
  \caption{DBMS\_C 综合执行链路图}
  \label{fig:integration-chain}
\end{figure}

这条链路可以按章节反向复盘：第8章 Parse 将 SQL 字符串变成 \codepath{QueryData} 或 \codepath{CommandData}；第9章 Plan/Scan 把 SELECT 变成计划树，再打开为运行期扫描链；第10章 Update 将更新命令分发到 \codepath{BasicUpdatePlanner}；第7章 Metadata 维护表、字段、视图和索引目录；第6章 Record 负责记录布局、slot 和表扫描；第5章 Transaction 连接写入、提交、回滚和并发锁基础；第4章 Buffer 维护页框缓存；第3章 Log 提供日志追加和刷盘基础；第2章 File/Page 完成文件块读写；第1章 Lib 为这些模块提供字符串、容器和字节缓冲等基础能力；第11章 Trace 用 CLI 文本观察这些过程。

表 \ref{tab:integration-command-module} 进一步把综合 SQL 与主要模块对应起来。实际执行时，一条命令通常会跨越多个模块，表格只列出最适合观察的主路径。

\begin{table}[htbp]
  \centering
  \footnotesize
  \setlength{\tabcolsep}{3pt}
  \caption{SQL 命令与模块关系表}
  \label{tab:integration-command-module}
  \begin{tabularx}{\textwidth}{>{\raggedright\arraybackslash}p{28mm}>{\raggedright\arraybackslash}p{50mm}Y}
    \toprule
    SQL 类型 & 主要经过模块 & 观察重点 \\
    \midrule
    CREATE TABLE & Parser、Update、Metadata、Record & 命令解析、表结构登记、系统目录写入 \\
    INSERT & Parser、Update、Record、Transaction、Buffer & 记录插入、事务写入、页框标脏 \\
    SELECT & Parser、Plan、Scan、Record、Metadata & QueryData、计划树、扫描链和输出行数 \\
    UPDATE & Parser、Update、Scan、Record、Transaction & 定位匹配记录并修改字段值 \\
    DELETE & Parser、Update、Scan、Record、Transaction & 定位匹配记录并删除 slot \\
    commit & Transaction、Log、Buffer & 提交事务、刷写日志和修改页 \\
    \bottomrule
  \end{tabularx}
\end{table}

\section{综合运行建议}

综合实验可以按两种方式观察。第一种是普通模式，重点观察 SQL 的用户可见输出；第二种是 Trace 模式，重点观察 SQL 在内部模块之间的传递过程。

\begin{dbcode}[listing options={language=sh}]{普通模式运行}
./bin/NewDBMS.exe
\end{dbcode}

普通模式适合确认命令是否可以正常执行、查询结果是否符合预期。Trace 默认关闭，因此普通输出不会混入 Trace 标签。

\begin{dbcode}[listing options={language=sh}]{Trace 模式运行}
./bin/NewDBMS.exe --trace
\end{dbcode}

也可以在交互式命令行中输入：

\begin{dbcode}[listing options={language=SQL}]{交互式开启 Trace}
trace on;
trace off;
trace;
\end{dbcode}

Trace 是 CLI 文本过程可视化，不是图形界面。它适合把第8章到第10章的执行链路连接起来观察，但不改变 SQL 的执行结果。表 \ref{tab:integration-normal-trace} 对普通模式和 Trace 模式作综合对照。

\begin{table}[htbp]
  \centering
  \small
  \setlength{\tabcolsep}{4pt}
  \caption{普通模式与 Trace 模式综合观察表}
  \label{tab:integration-normal-trace}
  \begin{tabularx}{\textwidth}{p{25mm}Y Y}
    \toprule
    运行方式 & 适合观察 & 注意事项 \\
    \midrule
    普通模式 & SQL 执行结果、查询表头和记录行 & 输出更接近用户使用场景，内部链路不可见 \\
    Trace 模式 & Parse、Plan、Scan、Update、Transaction、Output 等阶段 & Trace 默认关闭，需要通过 \codepath{--trace} 或 \codepath{trace on;} 开启 \\
    人工复盘 & 将输出与章节源码、图表和测试对应起来 & 建议使用新表名，避免旧数据库恢复影响演示现象 \\
    \bottomrule
  \end{tabularx}
\end{table}

\section{测试体系总览}

本指导书从第1章到第10章采用“先补真实 CTest，再写章节”的方式组织。第11章 Trace 是人工验证型实验，因此不新增自动测试。表 \ref{tab:integration-ctest-matrix} 汇总了已经作为章节依据的测试。

\begin{table}[htbp]
  \centering
  \scriptsize
  \setlength{\tabcolsep}{2.5pt}
  \caption{新增 CTest 覆盖矩阵}
  \label{tab:integration-ctest-matrix}
  \begin{tabularx}{\textwidth}{>{\raggedright\arraybackslash}p{38mm}p{15mm}p{31mm}Y}
    \toprule
    测试名称 & 对应章节 & 覆盖模块 & 验证重点 \\
    \midrule
    LibBasicTest & 第1章 & Lib & 字符串、容器、字节缓冲和词法切分基础行为 \\
    FileManagerBasicTest & 第2章 & File/Page & Page 读写、BlockID、FileManager 块级读写 \\
    LogManagerBasicTest & 第3章 & Log & 日志追加、刷盘和基础读取能力 \\
    BufferManagerBasicTest & 第4章 & buffer & pin/unpin、modified/dirty 和 flush 行为 \\
    TransactionBasicTest & 第5章 & tx & Transaction 与 BufferList、BufferManager 的基础协作 \\
    ConcurrencyBasicTest & 第5章 & tx/concurrency & S 锁、X 锁和释放锁的基础调用 \\
    RecoveryBasicTest & 第5章 & tx/recovery & 修改旧值记录和 rollback 基础行为 \\
    RecordBasicTest & 第6章 & record & Schema、Layout、RID、RecordPage、TableScan 基础行为 \\
    MetadataBasicTest & 第7章 & metadata & 表、视图、索引元数据登记与读取 \\
    ParseBasicTest & 第8章 & parse/query & SELECT、INSERT、UPDATE、DELETE、CREATE 的解析结果 \\
    PlanScanBasicTest & 第9章 & plan/query/record & SELECT 计划生成、Scan 打开和结果读取 \\
    UpdateBasicTest & 第10章 & plan/metadata/record & CREATE、INSERT、UPDATE、DELETE、VIEW、INDEX 更新路径 \\
    \bottomrule
  \end{tabularx}
\end{table}

这些测试不是为了覆盖所有边界条件，而是为实验指导书建立可运行的最小证据链。读者在扩展实验时，可以继续沿用同一原则：先确认真实 API 和可运行测试，再把测试路径、观察点和源码位置写入文档。

\section{模块边界总结}

理解 DBMS\_C 的关键，不是把所有代码混在一起看，而是先看清各层边界。表 \ref{tab:integration-chapter-module} 给出各章节与模块职责的总览。

\begin{table}[htbp]
  \centering
  \small
  \setlength{\tabcolsep}{4pt}
  \caption{各章节对应模块总览表}
  \label{tab:integration-chapter-module}
  \begin{tabularx}{\textwidth}{p{18mm}p{33mm}Y}
    \toprule
    章节 & 模块 & 职责边界 \\
    \midrule
    第1章 & Lib & 提供基础字符串、容器、字节缓冲和词法工具，不理解数据库语义 \\
    第2章 & File/Page & 负责页与文件块读写，不理解表和记录 \\
    第3章 & Log & 负责日志追加、刷盘和读取基础，不直接执行 SQL \\
    第4章 & Buffer & 管理页框、pin/unpin、dirty/flush，不理解 SQL \\
    第5章 & Transaction & 组织写入、提交、回滚和并发控制基础 \\
    第6章 & Record & 理解记录布局、slot 和表扫描 \\
    第7章 & Metadata & 管理表结构、字段、视图、索引等系统目录 \\
    第8章 & Parse & 只把 SQL 文本解析为数据对象，不执行查询 \\
    第9章 & Plan/Scan & 处理 SELECT 查询计划和运行期扫描 \\
    第10章 & Update & 处理更新类命令的执行分发 \\
    第11章 & Trace & 只观察执行过程，不改变执行语义 \\
    \bottomrule
  \end{tabularx}
\end{table}

从模块边界看，\codepath{File/Page} 不理解“表”，\codepath{Buffer} 不理解 SQL，\codepath{Record} 开始理解记录布局，\codepath{Metadata} 负责管理表结构，\codepath{Parser} 只解析，\codepath{Plan/Scan} 只处理查询执行，\codepath{Update} 处理更新命令，\codepath{Transaction/Log/Recovery} 支撑写入链路，\codepath{Trace} 只是旁路观察。

\section{常见问题与排查}

综合实验中常见的问题往往不是单个模块错误，而是运行环境、数据库目录或构建目录造成的观察偏差。

\begin{itemize}
  \item 重复运行同一个表名时，可能读到旧数据或旧系统目录记录。建议人工演示时使用新的表名，或在独立工作目录运行。
  \item Trace 默认关闭。如果看不到 \codepath{TRACE_PARSE}、\codepath{TRACE_PLAN} 等标签，先确认是否使用 \codepath{--trace} 或输入了 \codepath{trace on;}。
  \item Windows 下构建建议使用 MinGW Makefiles，并按当前项目验证过的命令执行；不要把 Linux 风格命令直接照搬到 PowerShell。
  \item 临时构建目录如果出现 \codepath{.tmp} 或 Access denied 等半更新现象，可以换新的 \codepath{build-check} 类目录重新验证当前目标。
  \item 不建议手动同步 CMake 生成的 \codepath{.tmp} 文件。生成目录异常时，应先停止并报告真实状态。
  \item 不建议把 CLI Trace 自动化测试作为必须步骤。Trace 更适合人工观察，核心解析、计划、更新路径已经由前面章节的 CTest 支撑。
\end{itemize}

\section{后续扩展方向}

本指导书定位为本科课程实践和项目阅读材料，后续可以沿以下方向继续扩展：

\begin{itemize}
  \item 更系统的错误处理实验，把 DBError 与更多模块入口逐步接合。
  \item 更稳定的测试数据库隔离机制，降低重复运行测试和人工演示时的目录残留影响。
  \item 更完整的索引查询路径实验，进一步观察索引元数据如何进入查询计划。
  \item 更清晰的执行计划可视化，把当前 CLI Trace 转换为结构化输出或可视化数据源。
  \item 更统一的文档排版和图表风格，提升后续教学使用时的阅读体验。
\end{itemize}

这些方向都可以继续坚持本指导书使用的实验组织方式：先确认真实代码和测试，再撰写章节说明。

\section{本章小结}

本章将 DBMS\_C 的前十一章内容收束为一条综合执行链路：SQL 先被解析为数据对象，SELECT 经由 Plan/Scan 执行，更新命令经由 BasicUpdatePlanner 执行，记录和元数据向下依赖事务、缓冲区、日志和文件页式存储，Trace 则提供 CLI 形式的过程观察。

作为一个教学型 C 语言 DBMS 原型，DBMS\_C 的价值在于它把数据库系统中常见的抽象压缩到可阅读、可运行、可测试的代码规模内。读者完成本指导书后，应能够从 SQL 输入一路追踪到底层文件块读写，并理解各模块之间如何协作完成一次数据库操作。
